%% ============================================================
%%  Aula 6 — Gradientes de política II: variância, confiança e PPO
%%  Versão sucinta: roteiro do apresentador (poucas palavras,
%%  mesmas definições e figuras do aula06.tex)
%%  Adaptação em pt-BR do CS234 (Stanford, Emma Brunskill), Lecture 6
%%  Seção de gradientes avançados baseada nos slides de Joshua Achiam
%%  Compilar: latexmk -xelatex aula06-sucinta.tex   (duas passadas)
%% ============================================================
\documentclass[aspectratio=169, 11pt]{beamer}
%% .sty do projeto na raiz do repositório (../)
\makeatletter
\def\input@path{{./}{../}}
\makeatother


\usetheme{AprendizadoRL}      % ANTES do babel — fontspec precisa vir primeiro
\usepackage{babel}
\babelprovide[import, main]{brazilian}
\usepackage{booktabs}
\usepackage{pgfplots}
\pgfplotsset{compat=1.18}
\usepgfplotslibrary{fillbetween}

\title[Aula 6 — Gradientes de política II]{Aula 6: Gradientes de política II\\ variância, confiança e PPO}
\subtitle[Aprendizagem por Reforço]{Linha de base \textbullet\ Ator--crítico \textbullet\ Limites de desempenho \textbullet\ PPO}
\author[Prof. Marcelo K. Albertini]{Prof. Marcelo Keese Albertini \\
  \footnotesize Universidade Federal de Uberlândia \\
  \footnotesize adaptado de Emma Brunskill, CS234 (Stanford); seção avançada de Joshua Achiam}
\institute{Universidade Federal de Uberlândia}
\date{2026}

% ---- Nota discreta: perguntas ficam para o professor conduzir em aula
\newcommand{\paradiscussao}{\smallskip
  {\footnotesize\color{rlCinza}\itshape Pergunta para discussão conduzida em aula.}}

% ---- Macros locais da aula
\newcommand{\Adv}{A}
\newcommand{\Ahat}{\hat{A}}
\newcommand{\ghat}{\hat{g}}
\newcommand{\polnovo}{\pol'}
\newcommand{\dpol}{d^{\pol}}
\newcommand{\dpolnovo}{d^{\pol'}}
\newcommand{\KLd}[2]{D_{\mathrm{KL}}\!\left(#1\,\|\,#2\right)}
\DeclareMathOperator{\clipop}{clip}
\newcommand{\Lsub}[1]{L_{#1}}
\newcommand{\Var}{\operatorname{Var}}

% O estilo `caixa` existe no .sty do handout, mas não no tema Beamer.
\tikzset{caixa/.style={draw, thick, rounded corners=2.5pt, align=center,
                       font=\scriptsize, inner sep=3pt}}

\begin{document}

\begin{frame}[plain, noframenumbering]\titlepage\end{frame}

%% ============================================================
\begin{frame}{Onde estamos}
  \framesubtitle{Três aulas sobre busca direta no espaço de políticas}

  \begin{columns}[T]
    \begin{column}{0.32\textwidth}
      \begin{block}{Aula 5 --- feito}
        \small Razão de verossimilhança, teorema do gradiente de política,
        REINFORCE.
      \end{block}
    \end{column}
    \begin{column}{0.32\textwidth}
      \begin{block}{Aula 6 --- hoje}
        \small REINFORCE \emph{funciona}, mas mal. Domar a variância; dar
        passos que não destroem a política.
      \end{block}
    \end{column}
    \begin{column}{0.32\textwidth}
      \begin{block}{Aula 7 --- a seguir}
        \small GAE, melhoria monotônica, aprendizagem por imitação.
      \end{block}
    \end{column}
  \end{columns}

  \vfill
  \begin{ideiabox}
    Aula 5: \emph{para onde} andar. Hoje: \alert{com que precisão} conhecemos
    essa direção, e \alert{que passo} dar sem cair.
  \end{ideiabox}
\end{frame}

%% ============================================================
\begin{frame}{Plano de hoje}
  \begin{enumerate}
    \item \textbf{Linha de base} --- subtrair $b(s)$ não muda o gradiente
          esperado, mas encolhe a variância; melhor escolha de $b$.
    \item \textbf{Alvos alternativos ao retorno de Monte Carlo} --- de $G_t$
          ao crítico: ator--crítico e a função vantagem.
    \item \textbf{Problemas estruturais do gradiente de política} ---
          eficiência amostral, colapso de desempenho, distância em
          \emph{parâmetros} vs.\ \emph{políticas}.
    \item \textbf{Limites de desempenho relativo} --- lema da diferença de
          desempenho, $\Lsub{\pol}(\polnovo)$, divergência KL.
    \item \textbf{PPO} --- penalidade KL adaptativa e objetivo recortado.
  \end{enumerate}

  \vfill
  {\footnotesize\color{rlCinza}Itens 3--5 seguem os slides de gradientes de
  política avançados de Joshua Achiam.}
\end{frame}

%% ============================================================
\begin{frame}{Retomando a Aula 5}
  \begin{quizbox}{}
    Selecione tudo o que é verdadeiro a respeito de gradientes de política:
    \begin{enumerate}\setlength{\itemsep}{0ex}\footnotesize
      \item $\nabla_\theta V(\theta) = \E_{\pol_\theta}\!\left[\nabla_\theta \log \pol_\theta(s,a)\, Q^{\pol_\theta}(s,a)\right]$
      \item $\theta$ é \emph{sempre} incrementado na direção de $\nabla_\theta \ln \pol(s_t, a_t, \theta)$
      \item Pares $(s,a)$ com $Q$ estimado maior têm, em média, sua probabilidade aumentada
      \item Há garantia de convergência ao ótimo global da classe de políticas
    \end{enumerate}
    \paradiscussao
  \end{quizbox}

  \pause
  \begin{respbox}
    \textbf{1 e 3.} Item 2 é falso: a direção é a do escore
    \emph{multiplicado} por $Q$, e se inverte com retorno negativo. Item 4 é
    falso: só há garantia de \alert{ótimo local}.
  \end{respbox}
\end{frame}

%% ============================================================
\section{Reduzindo a variância: a linha de base}

\begin{frame}{REINFORCE: o ponto de partida}
  \begin{algbox}{REINFORCE (gradiente de política de Monte Carlo)}
    \small
    Inicialize $\theta$ arbitrariamente\\
    \textbf{para cada} episódio $\{s_1, a_1, r_2, \dots, s_{T-1}, a_{T-1}, r_T\} \sim \pol_\theta$ \textbf{faça}\\
    \hspace*{1.2em}\textbf{para} $t = 1$ até $T-1$ \textbf{faça}\\
    \hspace*{2.4em}$\theta \leftarrow \theta + \alpha \nabla_\theta \log \pol_\theta(s_t, a_t)\, G_t$\\
    \hspace*{1.2em}\textbf{fim}\\
    \textbf{fim} \qquad \textbf{retorne} $\theta$
  \end{algbox}

  \pause
  \vspace{-0.6em}
  {\small Estimador com $m$ trajetórias:}
  \vspace{-0.4em}
  \[\textstyle
    \nabla_\theta V(\theta) \;\approx\; \frac{1}{m}\sum_{i=1}^{m}
      \rec(\tau^{(i)}) \sum_{t=0}^{T-1} \nabla_\theta \log \pol_\theta\!\left(a_t^{(i)} \mid s_t^{(i)}\right)
  \]
  \vspace{-1.1em}

  \pause
  \begin{alertabox}
    \textbf{Não-viesado, porém muito ruidoso.} Certo em média não basta:
    variância alta → passo aponta para a direção errada na maioria dos lotes.
  \end{alertabox}
\end{frame}

%% ============================================================
\begin{frame}{De onde vem a variância?}
  \framesubtitle{Três fontes distintas, com três remédios distintos}

  \begin{columns}[T]
    \begin{column}{0.5\textwidth}
      \begin{enumerate}\small\setlength{\itemsep}{0.5em}
        \item \textbf{Amostragem da trajetória.} $\tau$ é sorteada; o mesmo
              $\theta$ produz retornos diferentes a cada execução.
        \item \textbf{Escala e deslocamento do retorno.} Todos
              $\approx 100$: o gradiente empurra todo ``para cima''; o sinal
              útil é a \emph{diferença} entre eles.
        \item \textbf{Crédito temporal.} $\rec(\tau)$ multiplica o escore de
              \alert{todas} as ações, inclusive as posteriores à recompensa
              --- ruído puro.
      \end{enumerate}
    \end{column}
    \begin{column}{0.48\textwidth}
      \begin{defbox}{Os três consertos}
        \small
        \begin{itemize}\setlength{\itemsep}{0.25ex}
          \item \textbf{Estrutura temporal:} troque $\rec(\tau)$ por $G_t$
                (só o futuro conta) --- visto na Aula 5.
          \item \textbf{Linha de base:} subtraia $b(s_t)$ --- \emph{hoje}.
          \item \textbf{Alvos alternativos:} troque $G_t$ por um crítico
                --- \emph{hoje}.
        \end{itemize}
      \end{defbox}
      \begin{ideiabox}
        Os dois primeiros \alert{gratuitos}: variância menor, sem viés. O
        terceiro troca viés por variância deliberadamente.
      \end{ideiabox}
    \end{column}
  \end{columns}
\end{frame}

%% ============================================================
\begin{frame}{A linha de base}
  \begin{defbox}{Gradiente de política com linha de base}
    \[
      \nabla_\theta \E_\tau[\rec] = \E_\tau\!\left[
        \sum_{t=0}^{T-1} \nabla_\theta \log \pol(a_t \mid s_t; \theta)
        \left(\sum_{t'=t}^{T-1} r_{t'} - \mdestaque{rlLaranja}{b(s_t)}\right)
      \right]
    \]
  \end{defbox}

  \vspace{-0.2em}
  \begin{itemize}\small
    \item \alert{Qualquer} $b$ que dependa só do estado, e não da ação:
          segue não-viesado.
    \item Quase ótima: o retorno esperado,
          $b(s_t) \approx \E[r_t + \cdots + r_{T-1}]$.
  \end{itemize}

  \begin{ideiabox}
    \textbf{Leitura.} Aumente a log-probabilidade de $a_t$ na proporção em que
    o retorno foi \emph{melhor do que o esperado} naquele estado. Sem linha de
    base, ação medíocre em estado excelente é reforçada como excelente.
  \end{ideiabox}
\end{frame}

%% ============================================================
\begin{frame}{Por que isso importa: um exemplo de três ações}
  \framesubtitle{Mesmo estado, três ações, retornos 98, 101 e 104}

  \begin{columns}[T]
    \begin{column}{0.49\textwidth}
      \centering
      {\footnotesize\color{rlCinza}\bfseries Sem linha de base --- pesos $G_t$}\\[0.2em]
      \begin{tikzpicture}[scale=0.55]
        \draw[rlCinza!40] (0,-0.25) -- (0,3.1);
        \foreach \i/\g/\lab/\ac in {0/2.45/98/1, 1/2.53/101/2, 2/2.60/104/3}{
          \draw[trans destac, line width=2pt] (\i*1.35+0.55,0) -- (\i*1.35+0.55,\g);
          \node[font=\scriptsize, color=rlCinza] at (\i*1.35+0.55,-0.35) {$a_{\ac}$};
          \node[font=\scriptsize, color=rlLaranja] at (\i*1.35+0.55,\g+0.28) {\lab};}
        \draw[rlCinza!40] (0,0) -- (4.3,0);
      \end{tikzpicture}\\[0.4em]
      {\scriptsize Sinal ($\pm 3$) afogado no nível ($\approx 100$).}
    \end{column}
    \begin{column}{0.49\textwidth}
      \centering
      {\footnotesize\color{rlCinza}\bfseries Com $b(s)=101$ --- pesos $G_t - b(s)$}\\[0.2em]
      \begin{tikzpicture}[scale=0.55]
        \draw[rlCinza!40] (0,-1.6) -- (0,1.6);
        \draw[rlCinza!40] (0,0) -- (4.3,0);
        \draw[trans, rlVermelho, line width=2pt] (0.55,0) -- (0.55,-1.35);
        \node[font=\scriptsize, color=rlVermelho] at (0.55,-1.62) {$-3$};
        \node[font=\scriptsize, color=rlCinza] at (1.9,0.3) {$0$};
        \fill[rlCinza] (1.9,0) circle (1.6pt);
        \draw[trans destac, line width=2pt] (3.25,0) -- (3.25,1.35);
        \node[font=\scriptsize, color=rlLaranja] at (3.25,1.62) {$+3$};
        \node[font=\scriptsize, color=rlCinza] at (0.55,0.3) {$a_1$};
        \node[font=\scriptsize, color=rlCinza] at (3.25,-0.3) {$a_3$};
        \node[font=\scriptsize, color=rlCinza] at (1.9,-0.3) {$a_2$};
      \end{tikzpicture}\\[0.4em]
      {\scriptsize $a_1$ \emph{desencorajada}, $a_3$ encorajada. Média do passo
      igual; ruído muito menor.}
    \end{column}
  \end{columns}

  \vspace{0.4em}
  \begin{ideiabox}
    O \alert{gradiente esperado é idêntico} nos dois painéis. Muda quantas
    trajetórias a média amostral precisa para chegar perto dele.
  \end{ideiabox}
\end{frame}

%% ============================================================
\begin{frame}{A linha de base não introduz viés}
  \framesubtitle{Basta mostrar que o termo extra tem esperança nula}

  \vspace{-0.3em}
  {\footnotesize
  \begin{align*}
    \E_\tau&\left[\nabla_\theta \log \pol(a_t \mid s_t;\theta)\, b(s_t)\right] \\[0.1em]
    &= \E_{s_{0:t},\, a_{0:(t-1)}}\Big[\E_{s_{(t+1):T},\, a_{t:(T-1)}}
        \left[\nabla_\theta \log \pol(a_t \mid s_t;\theta)\, b(s_t)\right]\Big]
        &&\text{\scriptsize(quebra da esperança)}\\
    &= \E_{s_{0:t},\, a_{0:(t-1)}}\Big[b(s_t)\,\E_{a_t}
        \left[\nabla_\theta \log \pol(a_t \mid s_t;\theta)\right]\Big]
        &&\text{\scriptsize($b$ sai da esperança interna)}\\
    &= \E_{s_{0:t},\, a_{0:(t-1)}}\left[b(s_t)\sum_a \pol_\theta(a \mid s_t)
        \frac{\nabla_\theta \pol(a \mid s_t;\theta)}{\pol_\theta(a \mid s_t)}\right]
        &&\text{\scriptsize(razão de verossimilhança, ao contrário)}\\
    &= \E_{s_{0:t},\, a_{0:(t-1)}}\left[b(s_t)\, \nabla_\theta \sum_a \pol(a \mid s_t;\theta)\right]
      \;=\; \E_{s_{0:t},\, a_{0:(t-1)}}\left[b(s_t)\, \nabla_\theta 1\right] \;=\; 0
        &&\text{\scriptsize(normalização)}
  \end{align*}}

  \vspace{-1.2em}
  \begin{ideiabox}
    Tudo se apoia em um fato só: $\sum_a \pol(a\mid s;\theta) = 1$
    \alert{para todo $\theta$}, logo o gradiente dessa constante é nulo. É por
    isso que $b$ \emph{não pode} depender de $a$.
  \end{ideiabox}
\end{frame}

%% ============================================================
\begin{frame}{Por que a linha de base reduz a variância}
  \vspace{-0.4em}
  {\small
  A variância se decompõe por passo de tempo; concentre-se em um termo. Como a
  esperança \alert{não depende de $b$} (acabamos de provar), minimizar a
  variância é minimizar só o momento de segunda ordem:
  \vspace{-0.3em}
  \[
    \argmin_b\; \E_{s\sim\dpol}\E_{a\sim\pol,\,G}\!\left[
      \left(\nabla_\theta \log \pol(a \mid s;\theta)\right)^2
      \left(G_t(s) - b(s)\right)^2\right]
  \]}
  \vspace{-1.2em}

  \pause
  \begin{teobox}{Linha de base de variância mínima}
    {\small
    É um problema de \alert{mínimos quadrados ponderados}. Derivando e igualando a zero:
    \vspace{-0.3em}
    \[
      b^\star(s) =
      \frac{\E\!\left[\left(\nabla_\theta \log \pol(a\mid s;\theta)\right)^2 G_t(s)\right]}
           {\E\!\left[\left(\nabla_\theta \log \pol(a\mid s;\theta)\right)^2\right]}
      \;\approx\; \E\!\left[G_t(s)\right] = \Vpi(s)
    \]
    \vspace{-1.1em}
    \par\smallskip
    Isto é: a média dos retornos \emph{ponderada pelo quadrado do escore}.
    Ignorando essa ponderação, sobra a média simples --- a \alert{função valor
    de estado}.}
  \end{teobox}
\end{frame}

%% ============================================================
\begin{frame}{Algoritmo de gradiente de política ``baunilha''}
  \begin{algbox}{Gradiente de política com linha de base aprendida}
    \small
    Inicialize o parâmetro da política $\theta$ e a linha de base $b$\\[0.25em]
    \textbf{para} iteração $= 1, 2, \dots$ \textbf{faça}\\
    \hspace*{1.2em} Colete um conjunto de trajetórias executando a política atual\\
    \hspace*{1.2em} Em cada passo $t$ de cada $\tau^i$: retorno
                    $G_t^i = \sum_{t'\ge t} r_{t'}^i$ e vantagem
                    $\Ahat_t^i = G_t^i - b(s_t^i)$\\
    \hspace*{1.2em} Reajuste a linha de base minimizando
                    $\sum_i \sum_t \left(b(s_t^i) - G_t^i\right)^2$\\
    \hspace*{1.2em} Atualize a política com $\ghat$, soma de termos
                    $\nabla_\theta \log \pol(a_t \mid s_t, \theta)\,\Ahat_t$
                    {\color{rlCinza}(via SGD ou Adam)}\\
    \textbf{fim}
  \end{algbox}

  \vspace{-0.2em}
  \begin{ideiabox}
    Método de \alert{dois modelos} (política + valor): a um passo do
    ator--crítico. $b$ ajustada com a \emph{mesma} amostra do gradiente:
    prático, mas correlaciona os dois.
  \end{ideiabox}
\end{frame}

%% ============================================================
\begin{frame}{Que outras linhas de base existem?}
  \begin{columns}[T]
    \begin{column}{0.52\textwidth}
      \begin{defbox}{Funções valor como linha de base}
        \small
        Função valor de estado--ação:
        \[ \Qpi(s,a) = \E_\pol\!\left[r_0 + \gamma r_1 + \gamma^2 r_2 + \cdots \mid s_0=s,\, a_0=a\right] \]
        Função valor de estado --- a linha de base natural:
        \[ \Vpi(s) = \E_\pol\!\left[r_0 + \gamma r_1 + \cdots \mid s_0=s\right]
                   = \E_{a\sim\pol}\!\left[\Qpi(s,a)\right] \]
      \end{defbox}
    \end{column}
    \begin{column}{0.46\textwidth}
      {\small Outras escolhas usadas na prática:}
      \begin{itemize}\small\setlength{\itemsep}{0.3ex}
        \item \textbf{Média do lote:} $b = \frac{1}{N}\sum G_t$. Constante,
              trivial, já ajuda muito.
        \item \textbf{Regressão sobre atributos:} $b(s) = w^\top x(s)$.
        \item \textbf{Rede neural} treinada por regressão sobre $G_t$ --- o
              \emph{crítico}.
      \end{itemize}
      \vspace{0.3em}
      \begin{alertabox}
        \small Normalizar as vantagens do lote
        ($\Ahat \leftarrow (\Ahat - \bar{\Ahat})/\sigma_{\Ahat}$): padrão em
        implementações --- \emph{introduz} viés pequeno.
      \end{alertabox}
    \end{column}
  \end{columns}
\end{frame}

%% ============================================================
\section{Alvos alternativos: do retorno ao crítico}

\begin{frame}{Escolhendo o alvo}
  \framesubtitle{$G_t$ é uma estimativa de $\Vpi(s_t)$ obtida de uma única execução}

  \begin{columns}[T]
    \begin{column}{0.5\textwidth}
      \begin{itemize}\small
        \item $G_t^i$: não-viesado para $\Vpi(s_t)$ --- mas é amostra
              \emph{única} de uma soma de muitas variáveis aleatórias.
        \item Variância menor \alert{introduzindo viés}, via
              \emph{bootstrapping} e aproximação de função.
        \item Mesmo dilema de TD versus MC (Aula 3), agora no alvo do
              gradiente de política.
      \end{itemize}
    \end{column}
    \begin{column}{0.48\textwidth}
      \begin{defbox}{Métodos ator--crítico}
        \small
        Mantêm representações explícitas \alert{da política} (o ator) e
        \alert{da função valor} (o crítico), e atualizam as duas.
        A estimativa de $V$ ou $Q$ é feita pelo crítico.
      \end{defbox}
      \begin{ideiabox}
        \small
        O ator age; o crítico julga. A3C
        {\color{rlCinza}(Mnih et al., ICML 2016)} popularizou a receita.
      \end{ideiabox}
    \end{column}
  \end{columns}
\end{frame}

%% ============================================================
\begin{frame}{O laço ator--crítico}
  \centering
  \begin{tikzpicture}[node distance=8mm, scale=0.86, every node/.style={transform shape}]
    \node[caixa, draw=rlLaranja, fill=rlLaranja!10, minimum width=26mm, minimum height=11mm,
          font=\small, align=center] (ator) at (0,1.6)
          {\textbf{Ator} \\[-0.2em] {\scriptsize política $\pol_\theta(a\mid s)$}};
    \node[caixa, draw=rlTeal, fill=rlTeal!10, minimum width=26mm, minimum height=11mm,
          font=\small, align=center] (critico) at (0,-1.6)
          {\textbf{Crítico} \\[-0.2em] {\scriptsize valor $V_w(s)$}};
    \node[caixa, draw=rlAzul, fill=rlAzul!8, minimum width=26mm, minimum height=11mm,
          font=\small, align=center] (amb) at (6.4,0)
          {\textbf{Ambiente} \\[-0.2em] {\scriptsize $P(s'\mid s,a)$, $\rec$}};

    \draw[trans destac] (ator.east) to[bend left=12]
      node[above, font=\scriptsize, pos=0.45] {ação $a_t$} (amb.north west);
    \draw[trans, rlAzulClaro] (amb.south west) to[bend left=12]
      node[below, font=\scriptsize, pos=0.55] {$s_{t+1},\; r_t$} (critico.east);
    \draw[trans, rlTeal, very thick] (critico.north) -- node[right, font=\scriptsize, align=left]
      {vantagem \\[-0.2em] $\Ahat_t = r_t + \gamma V_w(s_{t+1}) - V_w(s_t)$} (ator.south);
    \draw[trans, rlCinza, dashed] (amb.north) to[out=100, in=20, looseness=0.75]
      node[above, font=\scriptsize, pos=0.55] {estado $s_t$} (ator.north east);
  \end{tikzpicture}

  \vspace{0.3em}
  \begin{alertabox}
    O crítico transforma um retorno \emph{observado até o fim do episódio} em
    um sinal disponível \alert{no passo seguinte} --- mas, errado, o ator
    otimiza a coisa errada. Os dois aprendem juntos; a estabilidade dessa
    dança é metade da literatura da área.
  \end{alertabox}
\end{frame}

%% ============================================================
\begin{frame}{Do retorno à vantagem}
  \vspace{-0.3em}\small
  Partindo da forma com linha de base:
  \[
    \nabla_\theta \E_\tau[\rec] = \E_\tau\!\left[\sum_{t=0}^{T-1}
      \nabla_\theta \log \pol(a_t\mid s_t;\theta)\left(\sum_{t'=t}^{T-1} r_{t'} - b(s_t)\right)\right]
  \]
  \pause\vspace{-0.5em}
  Trocando o retorno amostral por um crítico $Q(s_t,a_t;w)$:
  \vspace{-0.3em}
  \[
    \nabla_\theta \E_\tau[\rec] \approx \E_\tau\!\left[\sum_{t=0}^{T-1}
      \nabla_\theta \log \pol(a_t\mid s_t;\theta)\left(Q(s_t,a_t;w) - b(s_t)\right)\right]
  \]
  \pause\vspace{-0.5em}
  Com $b$ = estimativa de $V$, colapsa na \alert{função vantagem}:
  \vspace{-0.3em}
  \[
    \nabla_\theta \E_\tau[\rec] \approx \E_\tau\!\left[\sum_{t=0}^{T-1}
      \nabla_\theta \log \pol(a_t\mid s_t;\theta)\,
      \mdestaque{rlLaranja}{\Ahat^{\pol}(s_t,a_t)}\right]
  \]

  \vspace{-0.5em}
  \begin{defbox}{Função vantagem}
    \small $\Adv^{\pol}(s,a) = \Qpi(s,a) - \Vpi(s)$: quanto $a$ é melhor do que
    a ação \emph{média} de $\pol$ em $s$. Note que
    $\E_{a\sim\pol}[\Adv^{\pol}(s,a)] = 0$ para todo $s$.
  \end{defbox}
\end{frame}

%% ============================================================
\begin{frame}{O catálogo de alvos}
  \framesubtitle{Todos dão o mesmo gradiente esperado --- exceto os que não dão}
  \vspace{-0.3em}
  \centering\footnotesize
  \begin{tabular}{@{}llcc@{}}
    \toprule
    \textbf{Peso do escore} & \textbf{Nome} & \textbf{Viés} & \textbf{Variância} \\
    \midrule
    $\rec(\tau)$ & retorno total da trajetória & nenhum & altíssima \\
    $G_t = \sum_{t'\ge t} r_{t'}$ & retorno futuro (estrutura temporal) & nenhum & muito alta \\[-0.1em]
    $G_t - b(s_t)$ & retorno com linha de base & nenhum & alta \\
    $G_t - V_w(s_t)$ & vantagem de Monte Carlo & pequeno$^\dagger$ & média \\
    $\delta_t = r_t + \gamma V_w(s_{t+1}) - V_w(s_t)$ & erro de TD & maior & baixa \\
    $\Ahat_t^{\mathrm{GAE}(\gamma,\lambda)}$ & vantagem generalizada & ajustável & ajustável \\
    \bottomrule
  \end{tabular}

  \smallskip
  {\scriptsize $^\dagger$ o viés vem apenas do erro de aproximação de $V_w$;
  se $V_w = \Vpi$, não há viés.}

  \begin{ideiabox}
    Descer a tabela: do \emph{Monte Carlo} ao \emph{TD}. A última linha
    interpola entre as duas com um só $\lambda$ --- próxima aula.
  \end{ideiabox}
\end{frame}

%% ============================================================
\begin{frame}{Teste de compreensão --- linha de base e vantagem}
  \begin{quizbox}{}
    Trocamos o peso $G_t - b(s_t)$ por
    $\delta_t = r_t + \gamma V_w(s_{t+1}) - V_w(s_t)$.
    \begin{enumerate}\setlength{\itemsep}{0ex}\footnotesize
      \item A troca preserva a ausência de viés do estimador.
      \item A troca reduz a variância.
      \item Se $V_w = \Vpi$, então $\E[\delta_t \mid s_t, a_t] = \Adv^{\pol}(s_t,a_t)$.
      \item Uma linha de base $b(s,a)$ também manteria o estimador não-viesado.
    \end{enumerate}
    \paradiscussao
  \end{quizbox}
  \vspace{-0.5em}

  \pause
  \begin{respbox}
    \textbf{2 e 3.} Item 1 é falso: há viés assim que $V_w \ne \Vpi$ --- o
    preço que pagamos. Item 4 é falso: a prova de não-viés exige tirar $b$ de
    dentro de $\sum_a$.
  \end{respbox}
\end{frame}

%% ============================================================
\section{Os problemas do gradiente de política}

\begin{frame}{O problema de otimização, revisado}
  Os algoritmos de gradiente de política tentam resolver
  \[
    \max_\theta\; J(\pol_\theta) \defeq \E_{\tau\sim\pol_\theta}\!\left[\sum_{t=0}^{\infty}\gamma^t r_t\right]
  \]
  por subida de gradiente estocástica em $\theta$, usando
  \[
    g = \nabla_\theta J(\pol_\theta) = \E_{\tau\sim\pol_\theta}\!\left[
      \sum_{t=0}^{\infty}\gamma^t \nabla_\theta \log \pol_\theta(a_t\mid s_t)\,
      \Adv^{\pol_\theta}(s_t,a_t)\right]
  \]

  \pause
  \vspace{-0.4em}
  \begin{alertabox}
    \textbf{Três limitações estruturais:} \textbf{(1)} a eficiência amostral é
    ruim; \textbf{(2)} distância no \emph{espaço de parâmetros} $\ne$ distância
    no \emph{espaço de políticas}; \textbf{(3)} como consequência de (2),
    acertar o tamanho do passo é difícil --- e errar pode ser catastrófico.
  \end{alertabox}
\end{frame}

%% ============================================================
\begin{frame}{Limitação 1 --- eficiência amostral}
  \begin{columns}[T]
    \begin{column}{0.54\textwidth}
      \begin{itemize}
        \item Cada lote \alert{descartado imediatamente} após um único passo
              de gradiente.
        \item Gradiente de política = esperança \emph{sob a política atual}
              --- \textbf{on-policy} por construção. $\theta$ muda → dados
              envelhecem.
      \end{itemize}

      \medskip
      {\small Duas maneiras de obter uma estimativa não-viesada:}
      \begin{itemize}\small\setlength{\itemsep}{0.2ex}
        \item Amostrar trajetórias da própria política e formar a média
              amostral --- \emph{mais estável}.
        \item Usar trajetórias de outras políticas --- \emph{menos estável}.
      \end{itemize}
      \vspace{-0.4em}
    \end{column}
    \begin{column}{0.44\textwidth}
      \begin{ideiabox}
        \textbf{Oportunidade.} Reaproveitar dados antigos para dar
        \alert{vários} passos de gradiente antes de coletar novamente.
      \end{ideiabox}
      \begin{alertabox}
        \textbf{Desafio.} \emph{Quantos} passos podemos dar antes que os
        dados deixem de descrever a política que estamos otimizando?
      \end{alertabox}
      {\scriptsize\color{rlCinza}Pergunta que o PPO responde no fim da aula.}
    \end{column}
  \end{columns}
\end{frame}

%% ============================================================
\begin{frame}{Limitação 2 --- o tamanho do passo}
  \framesubtitle{$\theta_{k+1} = \theta_k + \alpha_k \ghat_k$: um passo ruim pode ser irrecuperável}

  \centering
  \begin{tikzpicture}
    \begin{axis}[
      width=10.8cm, height=3.9cm,
      axis lines=left, xlabel={\scriptsize parâmetros da política $\theta$},
      ylabel={\scriptsize desempenho $J(\pol_\theta)$},
      xlabel style={font=\scriptsize, color=rlCinza},
      ylabel style={font=\scriptsize, color=rlCinza},
      xtick=\empty, ytick=\empty,
      axis line style={rlCinza!60},
      domain=0:5.2, samples=220, clip=false, ymin=-4.6, ymax=3.8]
      \addplot[rlAzul, very thick, smooth]
        {3*exp(-(x-1.9)^2/0.75) - 4.2*exp(-(x-3.55)^2/0.16)};
      % ponto atual
      \fill[rlCinza] (axis cs:1.05,1.14) circle (2.4pt);
      \node[font=\scriptsize, color=rlCinza, anchor=east] at (axis cs:0.95,1.35) {$\theta_k$};
      % passo bom
      \draw[trans destac, line width=1.4pt] (axis cs:1.05,1.6) to[bend left=22] (axis cs:1.85,3.2);
      \node[font=\scriptsize, color=rlLaranja, anchor=south] at (axis cs:1.5,3.25) {passo adequado};
      % passo ruim
      \draw[trans, rlVermelho, line width=1.4pt] (axis cs:1.15,1.0) to[bend right=16] (axis cs:3.5,-3.7);
      \node[font=\scriptsize, color=rlVermelho, anchor=west] at (axis cs:3.62,-3.3) {colapso de desempenho};
    \end{axis}
  \end{tikzpicture}

  \vspace{-0.2em}
  \begin{alertabox}
    \textbf{Grande demais:} colapso --- \emph{auto-reforçado}, porque a
    política ruim coleta dados ruins. \textbf{Pequeno demais:} progresso
    lento. O tamanho ``certo'' \alert{muda com $\theta$}: nenhuma constante
    serve. Normalizar vantagens e Adam ajuda --- mas resolve?
  \end{alertabox}
\end{frame}

%% ============================================================
\begin{frame}{Limitação 3 --- o problema é mais do que o passo}
  \framesubtitle{Considere $\pol_\theta(a{=}1) = \sigma(\theta)$, \; $\pol_\theta(a{=}2) = 1-\sigma(\theta)$}

  \centering
  \begin{tikzpicture}
    \begin{axis}[
      width=10.2cm, height=3.5cm,
      axis lines=left, xlabel={\scriptsize parâmetro $\theta$},
      ylabel={\scriptsize $\pol_\theta(a{=}1)$},
      xlabel style={font=\scriptsize, color=rlCinza},
      ylabel style={font=\scriptsize, color=rlCinza},
      xtick={-1,0,1,4,5}, ytick={0,0.5,1},
      ticklabel style={font=\tiny, color=rlCinza},
      axis line style={rlCinza!60},
      domain=-7:7, samples=200, ymin=-0.08, ymax=1.14, clip=false]
      \addplot[rlAzul, very thick] {1/(1+exp(-x))};
      % intervalo 1: theta em [-1,1]
      \draw[rlLaranja!30, line width=0pt, fill=rlLaranja!14]
        (axis cs:-1,0) rectangle (axis cs:1,1.02);
      \draw[trans destac] (axis cs:-1,0.269) -- (axis cs:-1,0.731);
      \node[font=\scriptsize, color=rlLaranja, anchor=south] at (axis cs:0,1.03) {$\Delta\theta = 2 \;\Rightarrow\; \Delta\pol \approx 0{,}46$};
      % intervalo 2: theta em [4,6]
      \draw[rlTeal!30, line width=0pt, fill=rlTeal!14]
        (axis cs:4,0) rectangle (axis cs:6,1.02);
      \node[font=\scriptsize, color=rlTeal, anchor=south] at (axis cs:5,1.03) {$\Delta\theta = 2 \;\Rightarrow\; \Delta\pol \approx 0{,}016$};
    \end{axis}
  \end{tikzpicture}

  \vspace{-0.2em}
  \begin{columns}[T]
    \begin{column}{0.47\textwidth}
      \begin{defbox}{Espaço de políticas}
        \small No caso tabular, o conjunto
        $\Pi = \{\pol \in \R^{\abs{\gS}\times\abs{\gA}} :
           \sum_a \pol_{sa} = 1,\; \pol_{sa}\ge 0\}$.
      \end{defbox}
    \end{column}
    \begin{column}{0.51\textwidth}
      \begin{ideiabox}
        \small \textbf{A pergunta que organiza o resto da aula:} atualizar de
        modo a \alert{nunca mudar a política mais do que pretendíamos}?
      \end{ideiabox}
    \end{column}
  \end{columns}
\end{frame}

%% ============================================================
\section{Limites de desempenho relativo}

\begin{frame}{O que queremos de uma regra de atualização}
  \begin{itemize}\small
    \item Trajetórias da política \alert{mais recente} com a maior eficiência
          possível.
    \item Passos que respeitem distância no \alert{espaço de políticas}, não
          no espaço de parâmetros.
  \end{itemize}
  \vspace{-0.2em}
  Para isso, precisamos de uma identidade que relacione o desempenho de
  \emph{duas} políticas.

  \begin{teobox}{Lema da diferença de desempenho}
    {\small Para quaisquer políticas $\pol$ e $\polnovo$,
    \vspace{-0.3em}
    \[
      J(\polnovo) - J(\pol)
      = \E_{\tau\sim\polnovo}\!\left[\sum_{t=0}^{\infty}\gamma^t \Adv^{\pol}(s_t,a_t)\right]
      = \frac{1}{1-\gamma}\,\E_{\substack{s\sim\dpolnovo\\ a\sim\polnovo}}\!\left[\Adv^{\pol}(s,a)\right]
    \]}
  \end{teobox}

  \vspace{-0.3em}
  \begin{ideiabox}
    Leitura: \alert{a melhoria de $\polnovo$ sobre $\pol$ é a vantagem média,
    medida com a régua antiga $\Adv^{\pol}$, ao longo das trajetórias novas.}
    O resto da aula: tornar essa expressão computável.
  \end{ideiabox}
\end{frame}

%% ============================================================
\begin{frame}{A distribuição descontada de estados futuros}
  \vspace{-0.3em}
  \begin{defbox}{$\dpol$}
    \small
    $\dpol(s) = (1-\gamma)\sum_{t\ge 0}\gamma^t\, \Prob(s_t = s \mid \pol)$ ---
    a frequência com que $\pol$ visita $s$, com visitas tardias descontadas por
    $\gamma^t$. O fator $(1-\gamma)$ normaliza: $\sum_s \dpol(s) = 1$.
  \end{defbox}

  \pause
  \vspace{-0.5em}
  {\small Substituindo no lema; esperança sobre ações por \alert{amostragem
  por importância}:}
  \vspace{-0.4em}
  \[
    J(\polnovo) - J(\pol)
    = \frac{1}{1-\gamma}\E_{\substack{s\sim\dpolnovo\\ a\sim\polnovo}}\!\left[\Adv^{\pol}(s,a)\right]
    = \frac{1}{1-\gamma}\E_{\substack{s\sim\dpolnovo\\ a\sim\pol}}
      \!\left[\frac{\polnovo(a\mid s)}{\pol(a\mid s)}\, \Adv^{\pol}(s,a)\right]
  \]

  \vspace{-0.6em}

  \pause
  \begin{alertabox}
    Quase lá --- as \emph{ações} já vêm da política antiga. Resta $s \sim
    \dpolnovo$: os \alert{estados} ainda são visitados pela política nova, que
    é justamente a que ainda não temos.
  \end{alertabox}
\end{frame}

%% ============================================================
\begin{frame}{Uma aproximação útil}
  E se simplesmente disséssemos $\dpolnovo \approx \dpol$?
  \vspace{-0.2em}
  \[
    J(\polnovo) - J(\pol) \;\approx\; \frac{1}{1-\gamma}
      \E_{\substack{s\sim\dpol\\ a\sim\pol}}\!\left[
      \frac{\polnovo(a\mid s)}{\pol(a\mid s)}\,\Adv^{\pol}(s,a)\right]
    \;\defeq\; \Lsub{\pol}(\polnovo)
  \]

  \pause
  \vspace{-0.4em}
  \begin{teobox}{Limite de desempenho relativo \; {\normalfont\scriptsize (Achiam, Held, Tamar e Abbeel, 2017)}}
    {\small
    \vspace{-0.4em}
    \[
      \abs{\,J(\polnovo) - \left(J(\pol) + \Lsub{\pol}(\polnovo)\right)}
      \;\le\; C\,\sqrt{\E_{s\sim\dpol}\!\left[\KLd{\polnovo}{\pol}[s]\right]}
    \]}
  \end{teobox}

  \vspace{-0.3em}
  \begin{ideiabox}
    Erro de trocar $\dpolnovo$ por $\dpol$ controlado pela
    \alert{divergência KL média entre as políticas}. Próximas em KL → a
    aproximação é boa, e $\Lsub{\pol}(\polnovo)$ vira \emph{substituto}
    legítimo do objetivo verdadeiro.
  \end{ideiabox}
\end{frame}

%% ============================================================
\begin{frame}{Divergência de Kullback--Leibler}
  \begin{columns}[T]
    \begin{column}{0.5\textwidth}
      \begin{defbox}{$D_{\mathrm{KL}}$}
        \small Para distribuições $P$ e $Q$ sobre uma variável discreta,
        \vspace{-0.3em}
        \[ \KLd{P}{Q} = \sum_x P(x)\log\frac{P(x)}{Q(x)} \]
        \vspace{-1.1em}
        \par Entre políticas, em um estado $s$:
        \vspace{-0.3em}
        \[ \KLd{\polnovo}{\pol}[s] = \sum_{a\in\gA}\polnovo(a\mid s)\log\frac{\polnovo(a\mid s)}{\pol(a\mid s)} \]
      \end{defbox}
    \end{column}
    \begin{column}{0.48\textwidth}
      \textbf{Propriedades}
      \begin{itemize}\small\setlength{\itemsep}{0.2ex}
        \item $\KLd{P}{P} = 0$
        \item $\KLd{P}{Q} \ge 0$
        \item $\KLd{P}{Q} \ne \KLd{Q}{P}$ --- \alert{não é simétrica}, logo não
              é uma métrica.
      \end{itemize}
      \vspace{0.3em}
      \begin{ideiabox}
        \small Régua certa aqui: mede diferença \emph{entre distribuições}, e
        não entre os parâmetros que as geram --- a distinção do slide da
        sigmoide.
      \end{ideiabox}
    \end{column}
  \end{columns}
\end{frame}

%% ============================================================
\begin{frame}{O que ganhamos com a aproximação}
  \vspace{-0.2em}
  \[
    \Lsub{\pol}(\polnovo)
    = \frac{1}{1-\gamma}\E_{\substack{s\sim\dpol\\ a\sim\pol}}\!\left[
      \frac{\polnovo(a\mid s)}{\pol(a\mid s)}\Adv^{\pol}(s,a)\right]
    = \E_{\tau\sim\pol}\!\left[\sum_{t=0}^{\infty}\gamma^t
      \frac{\polnovo(a_t\mid s_t)}{\pol(a_t\mid s_t)}\Adv^{\pol}(s_t,a_t)\right]
  \]

  \begin{itemize}\small
    \item Lado direito estimável com trajetórias da política \alert{antiga}
          $\pol$: otimizar $\polnovo$ sem coletar dados novos --- \alert{vários
          passos de gradiente por lote}.
    \item Parente da amostragem por importância; o peso depende \emph{só do
          passo atual}, não do histórico inteiro. Sem produto de $t$ razões,
          os pesos não explodem nem desaparecem.
  \end{itemize}

  \vspace{-0.2em}
  \begin{alertabox}
    Licença condicional: $\Lsub{\pol}(\polnovo)$ só aproxima o objetivo
    verdadeiro \alert{enquanto $\polnovo$ ficar perto de $\pol$ em KL}. Sem
    restrição, otimizamos um modelo fora da sua região de validade.
  \end{alertabox}
\end{frame}

%% ============================================================
\begin{frame}{A região de confiança}
  \centering
  \begin{tikzpicture}[scale=0.92, every node/.style={transform shape}]
    % curvas de nivel do objetivo substituto
    \foreach \r/\o in {2.9/12, 2.2/22, 1.5/34, 0.8/50}{
      \draw[rlAzulClaro!\o!white, line width=1.1pt, rotate=-22]
        (1.9,0.55) ellipse [x radius=\r, y radius=\r*0.52];}
    \node[font=\scriptsize, color=rlAzulClaro] at (3.9,2.05) {curvas de nível de $\Lsub{\pol_k}$};
    % regiao de confianca
    \draw[rlLaranja, thick, dashed, fill=rlLaranja!8] (-0.6,-0.35) circle (1.15);
    \node[font=\scriptsize, color=rlLaranja, anchor=north] at (-0.6,-1.62)
      {$\big\{\polnovo : \bar{D}_{\mathrm{KL}}(\polnovo \| \pol_k) \le \delta\big\}$};
    % pontos
    \fill[rlAzul] (-0.6,-0.35) circle (2.4pt);
    \node[font=\scriptsize, color=rlAzul, anchor=east] at (-0.78,-0.35) {$\pol_k$};
    \fill[rlLaranja] (0.5,0.03) circle (2.4pt);
    \node[font=\scriptsize, color=rlLaranja, anchor=south west] at (0.55,0.05) {$\pol_{k+1}$};
    \draw[trans destac] (-0.6,-0.35) -- (0.45,0.0);
    % ponto ruim
    \fill[rlVermelho] (2.9,1.2) circle (2.2pt);
    \node[font=\scriptsize, color=rlVermelho, anchor=west] at (3.05,1.2)
      {máximo irrestrito de $\Lsub{\pol_k}$ --- fora da região de validade};
  \end{tikzpicture}

  \vspace{-0.2em}
  \begin{ideiabox}
    Maximizar o substituto \emph{dentro} da bola de KL: ideia do TRPO e, de
    forma aproximada e barata, do PPO. Fora da bola, o substituto deixa de
    informar sobre o objetivo real.
  \end{ideiabox}
\end{frame}

%% ============================================================
\begin{frame}{Teste de compreensão --- o substituto}
  \begin{quizbox}{}
    \vspace{-0.2em}Sobre $\Lsub{\pol}(\polnovo)$ e o limite de desempenho relativo:
    \begin{enumerate}\setlength{\itemsep}{0ex}\footnotesize
      \item $\Lsub{\pol}(\pol) = 0$.
      \item Em $\polnovo=\pol$, o gradiente de $\Lsub{\pol}$ coincide com o
            gradiente de política usual.
      \item O limite deixa de valer se $\polnovo$ estiver longe de $\pol$.
      \item Se $\Lsub{\pol}(\polnovo) > 0$, então $J(\polnovo) > J(\pol)$.
    \end{enumerate}
    \paradiscussao
  \end{quizbox}
  \vspace{-0.5em}

  \pause
  \begin{respbox}
    \textbf{1 e 2.} Em $\polnovo = \pol$ a razão vale $1$ e
    $\E_{a\sim\pol}[\Adv^{\pol}] = 0$; derivando, reobtemos o gradiente da
    Aula 5 --- $\Lsub{\pol}$ \emph{estende} aquele gradiente. Item 3 é falso:
    o limite vale sempre, mas a folga fica grande demais para informar --- daí
    o item 4 também ser falso.
  \end{respbox}
\end{frame}

%% ============================================================
\begin{frame}{Leituras recomendadas}
  \begin{itemize}
    \item \textbf{Kakade e Langford (2002)} --- \emph{Approximately Optimal
          Approximate Reinforcement Learning}. A origem da ideia de atualização
          conservadora de políticas.
    \item \textbf{Schulman et al. (2015)} --- \emph{Trust Region Policy
          Optimization} (TRPO). A restrição de KL imposta exatamente, via
          gradiente natural.
    \item \textbf{Achiam et al. (2017)} --- \emph{Constrained Policy
          Optimization}. Origem do limite apresentado neste slide.
    \item \textbf{Schulman et al. (2017)} --- \emph{Proximal Policy
          Optimization Algorithms}. O que veremos a seguir.
  \end{itemize}

  \vfill
  \begin{ideiabox}
    A sequência histórica é instrutiva: cada trabalho troca um pouco de rigor
    por muita simplicidade de implementação. O PPO é o extremo dessa troca ---
    e é o que todo mundo usa.
  \end{ideiabox}
\end{frame}

%% ============================================================
\section{Otimização proximal de política (PPO)}

\begin{frame}{PPO: duas maneiras de não se afastar demais}
  \framesubtitle{Uma família de métodos que impõe a restrição de KL \emph{aproximadamente}, sem gradiente natural}

  \begin{columns}[T]
    \begin{column}{0.49\textwidth}
      \begin{defbox}{Penalidade KL adaptativa}
        \small A atualização resolve um problema \alert{sem restrição}:
        \vspace{-0.3em}
        \[ \theta_{k+1} = \argmax_\theta\; \Lsub{\theta_k}(\theta)
           - \beta_k\, \bar{D}_{\mathrm{KL}}(\theta \| \theta_k) \]
        \vspace{-1.2em}
        \par com
        $\bar{D}_{\mathrm{KL}}(\theta\|\theta_k) =
         \E_{s\sim d^{\pol_k}}\!\left[\KLd{\pol_{\theta_k}(\cdot\mid s)}{\pol_\theta(\cdot\mid s)}\right]$.
        O coeficiente $\beta_k$ \alert{muda entre iterações} para fazer cumprir
        aproximadamente a restrição.
      \end{defbox}
    \end{column}
    \begin{column}{0.49\textwidth}
      \begin{defbox}{Objetivo recortado}
        \small Nenhuma penalidade explícita: a própria função objetivo é
        \emph{achatada} onde a política se afasta demais. É a variante mais
        usada --- e a que vocês implementarão.
      \end{defbox}
      \begin{ideiabox}
        \small TRPO resolve o problema com restrição de KL de forma exata e
        cara. O PPO faz um remendo barato que, empiricamente, funciona quase
        igual.
      \end{ideiabox}
    \end{column}
  \end{columns}
\end{frame}

%% ============================================================
\begin{frame}{PPO com penalidade KL adaptativa}
  \begin{algbox}{PPO --- penalidade KL adaptativa}
    \small
    \textbf{Entrada:} $\theta_0$, penalidade inicial $\beta_0$, KL alvo $\delta$\\[0.2em]
    \textbf{para} $k = 0,1,2,\dots$ \textbf{faça}\\
    \hspace*{1.2em} Colete trajetórias parciais $\mathcal{D}_k$ com $\pol_k = \pol(\theta_k)$\\
    \hspace*{1.2em} Estime as vantagens $\Ahat_t^{\pol_k}$ por qualquer método\\
    \hspace*{1.2em} Calcule
      $\theta_{k+1} = \argmax_\theta \Lsub{\theta_k}(\theta) - \beta_k \bar{D}_{\mathrm{KL}}(\theta\|\theta_k)$
      com $K$ passos de SGD por minilotes (Adam)\\
    \hspace*{1.2em} \textbf{se} $\bar{D}_{\mathrm{KL}}(\theta_{k+1}\|\theta_k) \ge 1{,}5\,\delta$
      \textbf{então} $\beta_{k+1} = 2\beta_k$\\
    \hspace*{1.2em} \textbf{senão se} $\bar{D}_{\mathrm{KL}}(\theta_{k+1}\|\theta_k) \le \delta/1{,}5$
      \textbf{então} $\beta_{k+1} = \beta_k/2$\\
    \textbf{fim}
  \end{algbox}

  \vspace{-0.2em}
  \begin{ideiabox}
    $\beta$ inicial \alert{importa pouco} --- adapta rápido. Algumas iterações
    violam a restrição; a maioria não. O \alert{$K$} é o ponto central: vários
    passos por lote, enfim.
  \end{ideiabox}
\end{frame}

%% ============================================================
\begin{frame}{PPO com objetivo recortado}
  \begin{defbox}{Razão de probabilidades e objetivo $L^{\mathrm{CLIP}}$}
    Seja $r_t(\theta) = \dfrac{\pol_\theta(a_t\mid s_t)}{\pol_{\theta_k}(a_t\mid s_t)}$. Então
    \vspace{-0.4em}
    \[
      L^{\mathrm{CLIP}}_{\theta_k}(\theta) = \E_{\tau\sim\pol_k}\!\left[
        \sum_{t=0}^{T}
        \min\!\Big(r_t(\theta)\,\Ahat_t^{\pol_k},\;
        \clipop\big(r_t(\theta),\, 1-\epsilon,\, 1+\epsilon\big)\,\Ahat_t^{\pol_k}\Big)
      \right]
    \]
    \vspace{-1.1em}
    \par onde $\epsilon$ é um hiperparâmetro (tipicamente $\epsilon = 0{,}2$).
    A atualização é $\theta_{k+1} = \argmax_\theta L^{\mathrm{CLIP}}_{\theta_k}(\theta)$.
  \end{defbox}

  \begin{itemize}\small
    \item $r_t(\theta_k) = 1$: no início de cada época de otimização, a razão
          vale exatamente $1$ e o objetivo coincide com $\Lsub{\theta_k}$.
    \item O $\min$ com a versão recortada torna o objetivo um
          \alert{limite inferior pessimista} do substituto --- nunca otimista
          quanto ao ganho de se afastar.
  \end{itemize}
\end{frame}

%% ============================================================
\begin{frame}{A geometria do recorte}
  \framesubtitle{$L^{\mathrm{CLIP}}$ em função de $r$, para um único passo de tempo}

  \centering
  \begin{columns}[T]
    \begin{column}{0.49\textwidth}
      \centering {\footnotesize\color{rlCinza}\bfseries vantagem $\Ahat > 0$}\\[-0.2em]
      \begin{tikzpicture}
        \begin{axis}[width=6.0cm, height=3.7cm, axis lines=middle,
          xlabel={\scriptsize $r$}, ylabel={\scriptsize $L^{\mathrm{CLIP}}$},
          xlabel style={font=\scriptsize, color=rlCinza, at={(ticklabel* cs:1)}, anchor=west},
          ylabel style={font=\scriptsize, color=rlCinza, at={(ticklabel* cs:1)}, anchor=south},
          xtick={0.8,1,1.2}, xticklabels={$1{-}\epsilon$,$1$,$1{+}\epsilon$},
          ytick=\empty, ticklabel style={font=\tiny, color=rlCinza},
          axis line style={rlCinza!60}, xmin=0, xmax=1.9, ymin=-0.15, ymax=1.65, clip=false]
          \addplot[rlAzul, very thick] coordinates {(0,0) (1.2,1.2) (1.9,1.2)};
          \draw[rlCinza!50, dashed] (axis cs:1.2,0) -- (axis cs:1.2,1.2);
          \fill[rlLaranja] (axis cs:1,1) circle (2pt);
          \node[font=\tiny, color=rlLaranja, anchor=north west] at (axis cs:1.02,0.98) {$r=1$};
        \end{axis}
      \end{tikzpicture}\\[-0.2em]
      {\scriptsize O ganho \emph{satura} em $1+\epsilon$: sem prêmio por
      aumentar mais a probabilidade da ação.}
    \end{column}
    \begin{column}{0.49\textwidth}
      \centering {\footnotesize\color{rlCinza}\bfseries vantagem $\Ahat < 0$}\\[-0.2em]
      \begin{tikzpicture}
        \begin{axis}[width=6.0cm, height=3.7cm, axis lines=middle,
          xlabel={\scriptsize $r$}, ylabel={\scriptsize $L^{\mathrm{CLIP}}$},
          xlabel style={font=\scriptsize, color=rlCinza, at={(ticklabel* cs:1)}, anchor=west},
          ylabel style={font=\scriptsize, color=rlCinza, at={(ticklabel* cs:0)}, anchor=north},
          xtick={0.8,1,1.2}, xticklabels={$1{-}\epsilon$,$1$,$1{+}\epsilon$},
          ytick=\empty, ticklabel style={font=\tiny, color=rlCinza},
          axis line style={rlCinza!60}, xmin=0, xmax=1.9, ymin=-1.65, ymax=0.15, clip=false]
          \addplot[rlAzul, very thick] coordinates {(0,-0.8) (0.8,-0.8) (1.9,-1.9)};
          \draw[rlCinza!50, dashed] (axis cs:0.8,0) -- (axis cs:0.8,-0.8);
          \fill[rlLaranja] (axis cs:1,-1) circle (2pt);
          \node[font=\tiny, color=rlLaranja, anchor=south west] at (axis cs:1.02,-1.0) {$r=1$};
        \end{axis}
      \end{tikzpicture}\\[-0.2em]
      {\scriptsize A punição \emph{não} satura à direita: ação ruim que ficou
      muito mais provável continua sendo penalizada.}
    \end{column}
  \end{columns}
\end{frame}

%% ============================================================
\begin{frame}{Por que o mínimo, e por que a assimetria}
  \begin{itemize}\small
    \item \textbf{$\Ahat_t > 0$ (ação boa).} O objetivo cresce com $r$ até
          $1+\epsilon$ e depois fica plano: o gradiente \alert{zera}. O
          otimizador para de empurrar a ação para cima quando ela já ficou
          $20\%$ mais provável.
    \item \textbf{$\Ahat_t < 0$ (ação ruim).} O objetivo é plano à esquerda de
          $1-\epsilon$ --- não vale a pena reduzir mais --- mas continua
          decrescendo à direita.
  \end{itemize}

  \vspace{-0.2em}
  \begin{ideiabox}
    A assimetria é deliberada. O $\min$ escolhe sempre o \alert{ramo mais
    pessimista}: o recorte remove o incentivo para se afastar, mas \emph{nunca}
    remove o incentivo para \emph{voltar}. Se um passo anterior empurrou uma
    ação ruim para $r \gg 1$, o gradiente ainda existe para corrigi-lo.
  \end{ideiabox}

  \vspace{-0.2em}
  \begin{alertabox}
    O recorte limita a \emph{razão}, não a divergência KL. Sem garantia formal
    de melhoria monotônica --- heurística barata inspirada na teoria da seção
    anterior.
  \end{alertabox}
\end{frame}

%% ============================================================
\begin{frame}{Teste de compreensão --- o recorte do PPO}
  \begin{quizbox}{}
    Considere os dois gráficos do slide anterior, com $\epsilon \in (0,1)$.
    \begin{enumerate}\setlength{\itemsep}{0ex}\footnotesize
      \item O gráfico da esquerda mostra $L^{\mathrm{CLIP}}$ quando
            $\Adv > 0$, e o da direita quando $\Adv < 0$.
      \item O gráfico da direita mostra $L^{\mathrm{CLIP}}$ quando
            $\Adv > 0$, e o da esquerda quando $\Adv < 0$.
      \item Depende do valor de $\epsilon$.
    \end{enumerate}
    \paradiscussao
  \end{quizbox}
  \vspace{-0.5em}

  \pause
  \begin{respbox}
    \textbf{1.} Com $\Adv > 0$ queremos $r$ grande, e o recorte satura o
    \emph{ganho} em $1+\epsilon$ --- patamar à direita. Com $\Adv < 0$ o
    objetivo é negativo e o patamar aparece à esquerda, em $1-\epsilon$. O
    valor de $\epsilon$ move os joelhos da curva, mas não troca os painéis.
  \end{respbox}
\end{frame}

%% ============================================================
\begin{frame}{PPO recortado: o algoritmo completo}
  \begin{algbox}{PPO com objetivo recortado}
    \small
    \textbf{Entrada:} $\theta_0$ (ator), $w_0$ (crítico), $\epsilon$, épocas $K$, tamanho do lote\\[0.2em]
    \textbf{para} $k = 0,1,2,\dots$ \textbf{faça}\\
    \hspace*{1.2em} Execute $\pol_{\theta_k}$ por $T$ passos em $N$ ambientes paralelos\\
    \hspace*{1.2em} Estime as vantagens $\Ahat_t$ (na prática, por GAE --- Aula 7)\\
    \hspace*{1.2em} Guarde $\log \pol_{\theta_k}(a_t\mid s_t)$: é o denominador de $r_t(\theta)$\\
    \hspace*{1.2em} \textbf{para} época $= 1,\dots,K$ \textbf{faça}\\
    \hspace*{2.4em} Para cada minilote: passo de Adam em
      $-L^{\mathrm{CLIP}} + c_1 L^{V} - c_2 \mathcal{H}[\pol_\theta]$\\
    \hspace*{1.2em} \textbf{fim}\\
    \textbf{fim}
  \end{algbox}

  \vspace{-0.2em}
  \begin{ideiabox}
    Três termos na perda: objetivo recortado do \alert{ator}; erro quadrático
    do \alert{crítico} ($L^V$); bônus de \alert{entropia} --- impede a política
    de virar determinística cedo demais.
  \end{ideiabox}
\end{frame}

%% ============================================================
\begin{frame}{PPO na prática}
  \begin{columns}[T]
    \begin{column}{0.5\textwidth}
      \textbf{Valores usuais}
      \begin{itemize}\small\setlength{\itemsep}{0.2ex}
        \item $\epsilon \approx 0{,}2$; \; $K \approx 3$--$10$ épocas por lote
        \item $c_1 \approx 0{,}5$ (crítico), $c_2 \approx 0{,}01$ (entropia)
        \item Vantagens normalizadas por minilote
        \item Corte de norma do gradiente em $0{,}5$
      \end{itemize}
      \vspace{0.3em}
      \begin{ideiabox}
        \small Monitore a KL empírica entre $\pol_\theta$ e $\pol_{\theta_k}$ e
        a \alert{fração recortada}. KL dispara → reduza $K$ ou a taxa de
        aprendizado.
      \end{ideiabox}
    \end{column}
    \begin{column}{0.48\textwidth}
      \begin{alertabox}
        \small \textbf{Armadilhas frequentes}
        \begin{itemize}\setlength{\itemsep}{0.2ex}
          \item Recalcular $\Ahat_t$ dentro do laço de épocas: as vantagens são
                de $\pol_{\theta_k}$ e devem ficar \emph{congeladas}.
          \item Esquecer de destacar o gradiente do denominador de $r_t$.
          \item $K$ grande demais: a política se afasta, e o recorte sozinho
                não segura a KL.
        \end{itemize}
      \end{alertabox}
      \vspace{0.2em}
      {\scriptsize\color{rlCinza}Boa parte do desempenho publicado do PPO vem
      dessas escolhas de implementação, e não do objetivo em si.}
    \end{column}
  \end{columns}
\end{frame}

%% ============================================================
\begin{frame}{O que você deve saber ao sair daqui}
  \begin{enumerate}\small\setlength{\itemsep}{0.35em}
    \item \textbf{Demonstrar} que uma linha de base $b(s)$ não introduz viés,
          e apontar exatamente onde a demonstração quebraria se $b$ dependesse
          da ação.
    \item \textbf{Explicar} por que a linha de base ótima é aproximadamente
          $\Vpi(s)$, e escrever o gradiente em termos da vantagem
          $\Adv^{\pol} = \Qpi - \Vpi$.
    \item \textbf{Distinguir} distância no espaço de parâmetros de distância no
          espaço de políticas, e dizer por que a segunda é a que importa.
    \item \textbf{Enunciar} o lema da diferença de desempenho e explicar qual
          aproximação transforma $\Lsub{\pol}(\polnovo)$ em algo computável com
          dados antigos --- e sob que condição ela é válida.
    \item \textbf{Escrever} $L^{\mathrm{CLIP}}$, desenhar seus dois ramos e
          justificar a assimetria entre eles.
  \end{enumerate}
\end{frame}

%% ============================================================
\begin{frame}{Resumo e próxima aula}
  \begin{columns}[T]
    \begin{column}{0.49\textwidth}
      \begin{defbox}{O fio da aula}
        \small
        \begin{itemize}\setlength{\itemsep}{0.2ex}
          \item Linha de base: variância menor, viés nenhum.
          \item Crítico: variância ainda menor, viés controlado.
          \item Substituto $\Lsub{\pol}$: vários passos por lote.
          \item Recorte: mantém o substituto dentro da sua região de validade.
        \end{itemize}
      \end{defbox}
    \end{column}
    \begin{column}{0.49\textwidth}
      \begin{ideiabox}
        \small \textbf{Aula 7}
        \begin{itemize}\setlength{\itemsep}{0.2ex}
          \item GAE: interpolar entre TD e MC com um parâmetro $\lambda$.
          \item Teoria de melhoria monotônica: quando o passo é
                \emph{garantidamente} bom.
          \item Aprendizagem por imitação: clonagem de comportamento, DAGGER,
                RL inversa de máxima entropia.
        \end{itemize}
      \end{ideiabox}
    \end{column}
  \end{columns}

  \vfill
  {\footnotesize\color{rlCinza}Material adaptado de \emph{CS234: Reinforcement
  Learning}, Emma Brunskill (Stanford); a seção de gradientes de política
  avançados segue os slides de Joshua Achiam.}
\end{frame}

\end{document}
